% s-176-p176-m002; original PDF 176
\Aphorism{s-176-p176-m002}{264}
Damit eine Landschaft uns zur Fruchtbarkeit stimme, muss sie nicht nur überhaupt stimmende Kraft haben, sondern ausserdem mit dem richtigen Zahn in unsere eigene Stimmung eingreifen — sowie man in gewissen Intervallen ziehen muss, um die Glocke in Schwung zu setzen. Andernfalls können die beiden Schwingungstendenzen, die von aussen erregte und die innerlich vorbereitete, einander zerstören: Interferenz zweier Wellensysteme, die einander vernichten oder verdoppeln, jenachdem Berg mit Thal oder Berg mit Berg zusammentrifft. Goethe sagt: hast in der bösen Stund geruht, ist dir die gute doppelt gut — die gute, wo Berg mit Berg, Thal mit Thal zusammentrifft.
\par

